% Generated draft: requires mathematical and source review.
\documentclass[11pt]{article}
\usepackage[margin=25mm]{geometry}
\usepackage{fontspec}
\setmainfont{Latin Modern Roman}
\usepackage{amsmath,amssymb,mathtools,unicode-math}
\setmathfont{Latin Modern Math}
\usepackage{enumitem,needspace,xurl,hyperref}
\setlength{\parindent}{0pt}
\setlength{\parskip}{5pt}
\setlength{\emergencystretch}{3em}
\newcommand{\N}{\mathbb{N}}
\newcommand{\Z}{\mathbb{Z}}
\newcommand{\Q}{\mathbb{Q}}
\newcommand{\R}{\mathbb{R}}
\newcommand{\C}{\mathbb{C}}
\newcommand{\F}{\mathbb{F}}
\newcommand{\A}{\mathbb{A}}
\newcommand{\PP}{\mathbb{P}}
\newcommand{\E}{\mathbb{E}}
\newcommand{\eps}{\varepsilon}
\newcommand{\dd}{\,\mathrm{d}}
\begin{document}
\section*{UnsolvedMath ID: 30002605}
\textbf{Source problem code:} catalogue entry
\subsection{Definitions and mathematical statement}
\textbf{Formulation status.} The supplied record does not specify a mathematical conjecture or a theorem to be proved. Its original question is: ``Why is peridynamics a useful mathematical model for mechanics, particularly as a multiscale model that behaves nonlocally at small scales, locally at large scales, and transitions between these regimes?''

Here peridynamics is the mechanics modeling framework named by the source. In descriptive terms, nonlocal behavior involves interactions between spatially separated points, whereas local behavior is described by quantities and derivatives at the same spatial point. These descriptions do not specify a particular peridynamic model or a limiting assertion.

\textbf{Unresolved mathematical specification.} No spatial domain, dimension, unknown fields, interaction law, scale parameter, governing equations, initial or boundary conditions, admissible data, or mode of convergence is supplied. In particular, the record does not specify whether local behavior at large scales means a limit as an interaction length tends to zero, a coarse-scale approximation, or another relation. Thus neither a quantified nonlocal-to-local statement nor a criterion for the usefulness or preference of peridynamics can be extracted without additional hypotheses.

The supplied statement-verification note identifies the question as a rhetorical heading followed by an explanation, rather than an unresolved mathematical problem, and assigns the statement status \texttt{unrecoverable}. The separate supplied \texttt{open} label does not establish a precise open conjecture. No replacement conjecture is asserted here.
\subsection{Short English statement}
The record asks why peridynamics is useful for mechanics across nonlocal and local scales. Its supplied verification identifies this as an explanatory heading, not a recoverable mathematical problem. A precise statement requires a specified model, scale regime, hypotheses, and mathematical target; these are absent.
\subsection{Sources}
\begin{itemize}
\item Local UnsolvedMath registry: \texttt{data/unsolvedmath/problems.json}, permanent ID \texttt{30002605}, source code \texttt{catalogue entry}, title ``Peridynamics as a Multiscale Mechanics Model,'' in the Oberwolfach Reports Open Problems collection. The supplied proposed year is 2014; no individual proposer is identified.
\item Supplied report citation: \emph{Computational Multiscale Methods} (2015), Oberwolfach Reports, pp.\ 1625--1681, DOI \texttt{10.4171/owr/2014/30}. \url{https://doi.org/10.4171/owr/2014/30}. The report is supplied as provenance for workshop discussions; its contents have not been independently consulted for this draft.
\item Supplied literature source: ``Oberwolfach Report 30/2014.'' \url{https://ems.press/content/serial-article-files/46520}.
\item Literature assessment dated \texttt{2026-08-21}: supplied dated curation, not independently verified. It labels the record open and describes modeling motivations involving nonlocal-to-local behavior and crack handling, while qualifying preference for peridynamics as application-dependent rather than a settled mathematical claim. The supplied background requests independent verification of time-sensitive status and cited claims.
\item Supplied statement-verification note reports a check against the official Oberwolfach report and identifies the original question as a rhetorical heading followed by an explanation of usefulness, with no unresolved mathematical problem stated. It records statement status \texttt{unrecoverable}, promotes no clean statement, and retains the original record for provenance. This verification is supplied metadata, not an independent verification performed for this draft.
\end{itemize}
\end{document}
